\PassOptionsToClass{onlytextwidth}{beamer}
%%
% This is an Overleaf template for presentations
% using the TUM Corporate Desing https://www.tum.de/cd
%
% For further details on how to use the template, take a look at our
% GitLab repository and browse through our test documents
% https://gitlab.lrz.de/latex4ei/tum-templates.
%
% The tumbeamer class is based on the beamer class.
% If you need further customization please consult the beamer class guide
% https://ctan.org/pkg/beamer.
% Additional class options are passed down to the base class.
%
% If you encounter any bugs or undesired behaviour, please raise an issue
% in our GitLab repository
% https://gitlab.lrz.de/latex4ei/tum-templates/issues
% and provide a description and minimal working example of your problem.
%%

%% ================================== base ================================== %%
\documentclass[
	german,            % define the document language (english, german)
	aspectratio=169,    % define the aspect ratio (169, 43)
	% handout=2on1,       % create handout with multiple slides (2on1, 4on1)
	% partpage=false,     % insert page at beginning of parts (true, false)
	% sectionpage=true,   % insert page at beginning of sections (true, false)
]{tumbeamer}

\usepackage{booktabs}
\usepackage{enumerate}
\usepackage{float}

% fvextra must be loaded before csquotes (minted pulls it in later)
\usepackage{fvextra}

%% ================================ language ================================ %%
\usepackage{babel}[ngerman]
\usepackage{csquotes}[autostyle]
\usepackage[useregional]{datetime2}
\usepackage{pmboxdraw}

%% ================================== math ================================== %%
\usepackage{amsmath}
\usepackage{amssymb}
\usepackage{ifthen}

%% ================================ graphics ================================ %%
\usepackage{graphicx}
\usepackage{xcolor}
\usepackage{pifont}
\usepackage{url}
\usepackage{hyperref}

\graphicspath{ {./resources/} }

%% ================================== tikz ================================== %%
\usepackage{pgf}
\usepackage{tikz}
\usepackage{pgfplots}
\usepackage{tikz-timing}
\usepackage{circuitikz}
\usepackage{karnaugh-map}

\usetikzlibrary{arrows, arrows.meta, backgrounds, positioning,
	shapes, shapes.geometric, shapes.symbols,
	patterns, patterns.meta,
	matrix, fit, calc,
	decorations.pathreplacing,
	tikzmark, chains,
	overlay-beamer-styles,
	automata}

% requires circuitikz >= 1.1.0
% for distros with older distributions, install TeX Live manually
% instead of using your package manager
% see: https://tug.org/texlive/quickinstall.html
\ctikzset{logic ports=european}

%% ============================= code listings ============================= %%
\usepackage[cache=true]{minted}
\usepackage{listings}

\usemintedstyle{borland}
\setminted{
	fontsize=\small,
	frame=none,
	breaklines=true,
}

%% ============================ colors & styling ============================ %%
\renewcommand{\thempfootnote}{\arabic{mpfootnote}}
\setbeamercolor{footnote}{fg=black}
\setbeamercolor{footnote mark}{fg=black}

%% ============================= custom macros ============================= %%
% https://tex.stackexchange.com/a/7045
\newcommand*\circled[1]{\tikz[baseline=(char.base)]{
		\node[shape=circle,draw,inner sep=2pt, font=\scriptsize] (char) {#1};}}

\newcommand*\colorcirc[1]{\tikz[baseline=(char.base)]{
		\node[shape=circle,fill=#1,inner sep=2pt] {};}}

\newcommand{\n}[1]{\overline{#1}}

%% ================================ metadata ================================ %%
\author{Niklas Ladurner}
\institute{\theChairName\\\theDepartmentName\\\theUniversityName}

\footline{\insertauthor~|~\insertshorttitle~|~\insertshortdate}

\TUMbeamersetup{
	title page = TUM tower,         % style of the title page
	part page = TUM toc,            % style of part pages
	section page = TUM toc,         % style of section pages
	content page = TUM more space,  % style of normal content pages
	tower scale = 1.0,              % scaling factor of TUM tower (if used)
	headline = TUM threeliner,      % which variation of headline to use
	footline = TUM default,         % which variation of footline to use
	% configure on which pages headlines and footlines should be printed
	headline on = {title page},
	footline on = {every page, title page=false},
}

% available frame styles for title page, part page, and section page:
% TUM default, TUM tower, TUM centered,
% TUM blue default, TUM blue tower, TUM blue centered,
% TUM shaded default, TUM shaded tower, TUM shaded centered,
% TUM flags
%
% additional frame styles for part page and section page:
% TUM toc
%
% available frame styles for content pages:
% TUM default, TUM more space
%
% available headline options:
% TUM empty, TUM oneliner, TUM twoliner, TUM threeliner, TUM logothreeliner
%
% available footline options:
% TUM empty, TUM default, TUM infoline


\title{Übung 01: Zahlensysteme}
\subtitle{Einführung in die Rechnerarchitektur}
\date{\DTMdisplaydate{2024}{10}{18}{-1}}

\begin{document}

\maketitle

\begin{frame}[c]{}{}
	\begin{center}
		\LARGE  Keine Garantie für die Richtigkeit der Tutorfolien.

		\Large Bei Unklarheiten/Unstimmigkeiten haben VL/ZÜ-Folien recht!
	\end{center}
\end{frame}

\begin{frame}[c]{Organisatorisches}{}
	\begin{itemize}
		\item Wer bin ich?
		\item Kommunikation:
		      \begin{itemize}
			      \item Zulip: ERA-Streams vs. Tutoriumsstream vs. DM
			      \item E-Mail
			      \item Fax (??)
		      \end{itemize}
		\item Mitschriften/Folien auf meiner Homepage: \url{https://home.in.tum.de/~ladu/}
		\item Anmerkungen zu den Hausaufgaben/Übungen
		\item Tutoriumszeiten
	\end{itemize}
\end{frame}

\begin{frame}[c]{Was macht so ein Rechner eigentlich?}{}
	\begin{itemize}
		\item Rechner arbeiten mit Zahlen: Bilder, Strings, ...
		\item Dezimalsystem ungeeignet: zu viele Zustände (Ziffern)
		\item Strom an bzw. aus mittels Transistoren
		\item zwei Zustände $\rightarrow$ Binärsystem
		\item Oktalsystem und Hexadezimalsystem passen gut dazu! \\ $\rightarrow$ Warum?
	\end{itemize}
\end{frame}

\begin{frame}[c]{Wie funktioniert denn so ein Zahlensystem?}{}

	\begin{equation}
		W=\sum_{i=0}^{n-1}a_i\cdot B^i
	\end{equation}
	\begin{block}{Beispiel.}
		\begin{enumerate}
			\item $(1010)_2 = 0\cdot 2^0+1\cdot 2^1+0\cdot 2^2+1\cdot 2^3=2+8=10$
			\item $(763)_8 = 3\cdot 8^0+6\cdot 8^1+7\cdot 8^2=3+48+448=499$
			\item $(123)_{16} = 3\cdot 16^0+2\cdot 16^1+1\cdot 16^2=3+32+256=291$
		\end{enumerate}
	\end{block}
	\hspace{1cm}
	\begin{itemize}
		\item Was passiert bei einem Verschieben der Stellen nach links/rechts?
		\item Wie kann man Binärzahlen in Hexadezimal/Oktal umwandeln?
	\end{itemize}
\end{frame}

\begin{frame}[c]{Rechnen im Binärsystem}{}
	\begin{itemize}
		\item schriftliches Rechnen grundsätzlich gleich wie im Dezimalsystem
		\item Subtraktion: Addition mit negativer Zahl
		\item Einerkomplement vs. Zweierkomplement. Datenwortbreite ausschlaggebend!
	\end{itemize}
	\vfill{}
	\begin{block}{Beispiel.}
		\begin{enumerate}
			\item $(1100.0110)_2 = -2^7+(2^1+2^2+2^6)=-58$ bei Breite von 8 Bit
			\item $(1010)_2 = -2^3 + (2^1) = -6$ bei Breite von 4 Bit
		\end{enumerate}
	\end{block}

\end{frame}

\begin{frame}[c]{}{}
	\begin{center}
		\LARGE Fragen?
	\end{center}
\end{frame}

\begin{frame}[c]{Artemis-Hausaufgaben}{}
	\begin{itemize}
		\item \enquote{H01 --- Arithmetik und negative Zahlen} bis 27.10.2024 23:59 Uhr
		\item 100\% bedeuten 100\%, keine \emph{hidden tests} nach der Deadline
		\item kurze Einführung in Artemis/Git
	\end{itemize}
\end{frame}

\begin{frame}[fragile, c]{Links}{}
	\begin{itemize}
		\item Zulip: \href{https://zulip.in.tum.de/#narrow/stream/2661-ERA-Tutorium---Do-1600-1}{\enquote{ERA Tutorium - Do-1600-1}}
		      bzw. \href{https://zulip.in.tum.de/#narrow/stream/2675-ERA-Tutorium---Fr-1500-2 }{\enquote{ERA Tutorium - Fr-1500-2}}
		\item \href{https://www.moodle.tum.de/course/view.php?id=100633}{ERA-Moodle-Kurs}
		\item \href{https://artemis.in.tum.de/courses/401}{ERA-Artemis-Kurs}
		\item \href{https://git-scm.com/docs/gittutorial}{Git-Tutorial}, \href{https://rogerdudler.github.io/git-guide/}{alternatives Tutorial}
	\end{itemize}
\end{frame}

\maketitle

\end{document}
